% GENERATED FILE. Edit canonical lessons, then run npm run generate:books.
% canonical-source-hash: fnv1a64:6bf2a0e35b462ddd
% canonical-lessons: IT-C320-fiera, IT-C320-palio, IT-C320-giostra, IT-C320-onomastico, IT-C320-santo, IT-R320-ripasso-313-320

\chapter{Feste di paese}
\label{ch:it-feste-di-paese}
\hlchaptermodality{320}

\begin{chapteropening}
\textbf{By the end of this chapter:} \emph{I can talk about town festivities: the fair, the horse race, the merry-go-round, name days and saints.}

Complete the last lesson of chapter 320: I can talk about town festivities: the fair, the horse race, the merry-go-round, name days and saints.
\end{chapteropening}

\section[la fiera]{la fiera — the fair, the trade show}

\label{lesson:IT-C320-fiera}



\begin{quote}
Before the new one: say the Italian for to dress up in costume, then the Italian for to parade.
\end{quote}

\subsection*{You'll want to know: la fiera}
\textbf{la fiera} — \textquotedblleft{}the fair, the trade show\textquotedblright{}.

\begin{quote}
A ottobre c'è la fiera del libro.
\end{quote}

\begin{quote}
\emph{In October there is the book fair.}
\end{quote}

An old, unrelated \emph{fiera} means a wild beast.

\begin{cousinweb}
\emph{fiera} \textquotedblleft{}fair\textquotedblright{} is inherited from Late Latin \emph{fēria} \textquotedblleft{}a holiday, a feast day\textquotedblright{} (Lewis \& Short, under \emph{fēriae} \textquotedblleft{}holidays\textquotedblright{}). EtymDB lists \emph{fiera} under Latin \emph{fera} \textquotedblleft{}wild beast\textquotedblright{}, the source of a different \emph{fiera}.
\end{cousinweb}

\begin{grammarlens}[title={What you've built}]
One more word for telling, explaining and describing at B1.
\end{grammarlens}

\subsection*{Your turn}
\begin{itemize}
\raggedright
  \item \emph{Say it:} \emph{la fiera}
  \item \emph{Say it:} \emph{la fiera}, once more
  \item \emph{Say it:} a sentence of your own with \emph{la fiera}
  \item \emph{Recall:} say the Italian for to dress up in costume, then the Italian for to parade, then say \emph{la fiera} again
\end{itemize}

\begin{tcolorbox}[breakable,colback=teal!4,colframe=teal!35!black,title={Before you move on}]
What does la fiera mean? (\textquotedblleft{}The fair, the trade show\textquotedblright{}.) Say it once more.
\end{tcolorbox}
\section[il palio]{il palio — the palio, the traditional horse race}

\label{lesson:IT-C320-palio}



\begin{quote}
Before the new one: say the Italian for to parade, then the Italian for the fair.
\end{quote}

\subsection*{You'll want to know: il palio}
\textbf{il palio} — \textquotedblleft{}the palio, the traditional horse race\textquotedblright{}.

\begin{quote}
Il palio si corre in piazza due volte l'anno.
\end{quote}

\begin{quote}
\emph{The palio is run in the square twice a year.}
\end{quote}

Named after the prize, a painted banner.

\begin{cousinweb}
\emph{palio} is inherited from Latin \emph{pallium} \textquotedblleft{}a cloak, a cloth\textquotedblright{} (Lewis \& Short): the race is run for a precious cloth.
\end{cousinweb}

\begin{grammarlens}[title={What you've built}]
One more word for telling, explaining and describing at B1.
\end{grammarlens}

\subsection*{Your turn}
\begin{itemize}
\raggedright
  \item \emph{Say it:} \emph{il palio}
  \item \emph{Say it:} \emph{il palio}, once more
  \item \emph{Say it:} a sentence of your own with \emph{il palio}
  \item \emph{Recall:} say the Italian for to parade, then the Italian for the fair, then say \emph{il palio} again
\end{itemize}

\begin{tcolorbox}[breakable,colback=teal!4,colframe=teal!35!black,title={Before you move on}]
What does il palio mean? (\textquotedblleft{}The palio, the traditional horse race\textquotedblright{}.) Say it once more.
\end{tcolorbox}
\section[la giostra]{la giostra — the merry-go-round, the funfair ride}

\label{lesson:IT-C320-giostra}



\begin{quote}
Before the new one: say the Italian for the fair, then the Italian for the palio.
\end{quote}

\subsection*{You'll want to know: la giostra}
\textbf{la giostra} — \textquotedblleft{}the merry-go-round, the funfair ride\textquotedblright{}.

\begin{quote}
I bambini vogliono andare sulla giostra.
\end{quote}

\begin{quote}
\emph{The children want to go on the merry-go-round.}
\end{quote}

Long ago, \emph{la giostra} was a joust between knights.

\begin{cousinweb}
\emph{giostra} was borrowed from Old French \emph{joste} \textquotedblleft{}a joust\textquotedblright{} (EtymDB), from Vulgar Latin \emph{iuxtāre} \textquotedblleft{}to come close\textquotedblright{}, built on \emph{iuxtā} \textquotedblleft{}near\textquotedblright{}.
\end{cousinweb}

\begin{grammarlens}[title={What you've built}]
One more word for telling, explaining and describing at B1.
\end{grammarlens}

\subsection*{Your turn}
\begin{itemize}
\raggedright
  \item \emph{Say it:} \emph{la giostra}
  \item \emph{Say it:} \emph{la giostra}, once more
  \item \emph{Say it:} a sentence of your own with \emph{la giostra}
  \item \emph{Recall:} say the Italian for the fair, then the Italian for the palio, then say \emph{la giostra} again
\end{itemize}

\begin{tcolorbox}[breakable,colback=teal!4,colframe=teal!35!black,title={Before you move on}]
What does la giostra mean? (\textquotedblleft{}The merry-go-round, the funfair ride\textquotedblright{}.) Say it once more.
\end{tcolorbox}
\section[l'onomastico]{l'onomastico — the name day}

\label{lesson:IT-C320-onomastico}



\begin{quote}
Before the new one: say the Italian for the palio, then the Italian for the merry-go-round.
\end{quote}

\subsection*{You'll want to know: l'onomastico}
\textbf{l'onomastico} — \textquotedblleft{}the name day\textquotedblright{}.

\begin{quote}
Oggi è il mio onomastico.
\end{quote}

\begin{quote}
\emph{Today is my name day.}
\end{quote}

The feast of the saint whose name you bear.

\begin{cousinweb}
\emph{onomastico} is usually traced to Greek \emph{onomastikós} \textquotedblleft{}of names\textquotedblright{}, built on \emph{ónoma} \textquotedblleft{}name\textquotedblright{} (EtymDB has no Italian entry).
\end{cousinweb}

\begin{grammarlens}[title={What you've built}]
One more word for telling, explaining and describing at B1.
\end{grammarlens}

\subsection*{Your turn}
\begin{itemize}
\raggedright
  \item \emph{Say it:} \emph{l'onomastico}
  \item \emph{Say it:} \emph{l'onomastico}, once more
  \item \emph{Say it:} a sentence of your own with \emph{l'onomastico}
  \item \emph{Recall:} say the Italian for the palio, then the Italian for the merry-go-round, then say \emph{l'onomastico} again
\end{itemize}

\begin{tcolorbox}[breakable,colback=teal!4,colframe=teal!35!black,title={Before you move on}]
What does l'onomastico mean? (\textquotedblleft{}The name day\textquotedblright{}.) Say it once more.
\end{tcolorbox}
\section[il santo]{il santo — the saint}

\label{lesson:IT-C320-santo}



\begin{quote}
Before the new one: say the Italian for the merry-go-round, then the Italian for the name day.
\end{quote}

\subsection*{You'll want to know: il santo}
\textbf{il santo} — \textquotedblleft{}the saint\textquotedblright{}.

\begin{quote}
Ogni paese festeggia il suo santo.
\end{quote}

\begin{quote}
\emph{Every village celebrates its own saint.}
\end{quote}

Before a man's name it usually becomes \emph{San}, or \emph{Sant'} before a vowel; a woman saint is \emph{santa}.

\begin{cousinweb}
\emph{santo} is inherited from Latin \emph{sānctus} \textquotedblleft{}holy\textquotedblright{} (EtymDB).
\end{cousinweb}

\begin{grammarlens}[title={What you've built}]
One more word for telling, explaining and describing at B1.
\end{grammarlens}

\subsection*{Your turn}
\begin{itemize}
\raggedright
  \item \emph{Say it:} \emph{il santo}
  \item \emph{Say it:} \emph{il santo}, once more
  \item \emph{Say it:} a sentence of your own with \emph{il santo}
  \item \emph{Recall:} say the Italian for the merry-go-round, then the Italian for the name day, then say \emph{il santo} again
\end{itemize}

\begin{tcolorbox}[breakable,colback=teal!4,colframe=teal!35!black,title={Before you move on}]
What does il santo mean? (\textquotedblleft{}The saint\textquotedblright{}.) Say it once more.
\end{tcolorbox}
\section[(dialogue)]{Ripasso dei capitoli 313–320}

\label{lesson:IT-R320-ripasso-313-320}



\begin{quote}
Say the Italian for death and the deceased.
\end{quote}

\subsection*{Your turn}
Say these aloud:

\begin{itemize}
\raggedright
  \item death, the deceased, the funeral, mourning, to bury
  \item the will, to inherit, the inheritance, the heir, the grave
  \item to celebrate formally, the recurring occasion, the anniversary, best wishes, to wish
  \item the birthday candle, the person being celebrated, to wrap, the bow, the balloon
  \item the sparkling wine, the cork, to uncork, the corkscrew, the banquet
  \item the invitation, the guest, the reception, the refreshments, the sugared almond
  \item the carnival, the mask, the confetti, to dress up in costume, to parade
  \item the fair, the palio, the merry-go-round, the name day, the saint
\end{itemize}

\begin{tcolorbox}[breakable,colback=teal!4,colframe=teal!35!black,title={Before you move on}]
Death? (\textbf{la morte}.) The deceased? (\textbf{il defunto}.) The funeral? (\textbf{il funerale}.) Mourning? (\textbf{il lutto}.) To bury? (\textbf{seppellire}.) The will? (\textbf{il testamento}.) To inherit? (\textbf{ereditare}.) The inheritance? (\textbf{l'eredità}.) The heir? (\textbf{l'erede}.) The grave? (\textbf{la tomba}.) To celebrate formally? (\textbf{celebrare}.) The recurring occasion? (\textbf{la ricorrenza}.) The anniversary? (\textbf{l'anniversario}.) Best wishes? (\textbf{gli auguri}.) To wish? (\textbf{augurare}.) The birthday candle? (\textbf{la candelina}.) The person being celebrated? (\textbf{il festeggiato}.) To wrap? (\textbf{incartare}.) The bow? (\textbf{il fiocco}.) The balloon? (\textbf{il palloncino}.) The sparkling wine? (\textbf{lo spumante}.) The cork? (\textbf{il tappo}.) To uncork? (\textbf{stappare}.) The corkscrew? (\textbf{il cavatappi}.) The banquet? (\textbf{il banchetto}.) The invitation? (\textbf{l'invito}.) The guest? (\textbf{l'invitato}.) The reception? (\textbf{il ricevimento}.) The refreshments? (\textbf{il rinfresco}.) The sugared almond? (\textbf{il confetto}.) The carnival? (\textbf{il carnevale}.) The mask? (\textbf{la maschera}.) The confetti? (\textbf{i coriandoli}.) To dress up in costume? (\textbf{mascherarsi}.) To parade? (\textbf{sfilare}.) The fair? (\textbf{la fiera}.) The palio? (\textbf{il palio}.) The merry-go-round? (\textbf{la giostra}.) The name day? (\textbf{l'onomastico}.) The saint? (\textbf{il santo}.)
\end{tcolorbox}
